% !TEX root = ../../ctfp-print.tex

\lettrine[lhang=0.17]{型}{と関数の圏は}プログラミングにおいて重要な役割を果たし
ている。だから、型とは何か、なぜそれが必要かについて話そう。

\section{誰が型を必要とするのか？}

静的型付けと動的型付け、強い型付けと弱い型付けの優劣については、いくつかの論争が
あるようだ。思考実験でこれらの選択肢を説明しよう。何百万匹ものサルがコンピュータ
のキーボードの前で嬉しそうにランダムなキーを叩き、プログラムを生成し、コンパイル
して実行している様子を想像してほしい。

\begin{figure}[H]
  \centering
  \includegraphics[width=0.3\textwidth]{images/img_1329.jpg}
\end{figure}

\noindent
機械語では、サルが生成したバイトの任意の組み合わせが受け入れられ実行される。しかし
高水準言語では、コンパイラが字句エラーや文法エラーを検出できるという事実を評価する。
多くのサルはバナナを得られないが、残ったプログラムは有用である可能性が高くなる。
型チェックはナンセンスなプログラムに対するさらなる障壁を提供する。さらに、動的型付け
言語では型の不一致は実行時に発見されるが、強く型付けされた静的チェック言語では型の
不一致はコンパイル時に発見され、多くの誤ったプログラムが実行される前に排除される。

では、サルを喜ばせたいのか、それとも正しいプログラムを生成したいのか？

タイプするサルの思考実験での通常の目標は、シェイクスピアの全集の生産だ。スペル
チェッカーと文法チェッカーをループに組み込むことで、確率は劇的に上がるだろう。
型チェッカーの類比はさらに進んで、一度ロミオが人間と宣言されたら、彼が葉を生やし
たり強力な重力場で光子を閉じ込めたりしないことを保証する。

\section{型は合成可能性についてのものだ}

圏論は射の合成についてのものだ。しかし任意の2つの射が合成できるわけではない。一方
の射の終域オブジェクトは、次の射の始域オブジェクトと同じでなければならない。プログ
ラミングでは、ある関数の結果を別の関数に渡す。もし受け取る関数がソース関数の生成した
データを正しく解釈できなければ、プログラムは動作しない。合成が機能するためには両端が
合っている必要がある。言語の型システムが強力であるほど、このマッチをより適切に記述し
機械的に検証できる。

強い静的型チェックに対して私が聞く唯一の真剣な論拠は、意味的に正しいプログラムの
一部を排除してしまうかもしれないというものだ。実際には、これは非常にまれにしか起こら
ないし、いずれにせよ、本当に必要なときに型システムをバイパスするための何らかの抜け穴
をすべての言語が提供している。Haskellでさえ\code{unsafeCoerce}がある。しかしそのよう
な手段は慎重に使うべきだ。フランツ・カフカの登場人物グレゴール・ザムザは、巨大な虫に
変身するときに型システムを壊す――その結末は誰もが知っている。

よく聞くもう一つの論拠は、型を扱うことがプログラマに過度な負担を課すというものだ。
C++でイテレータの宣言をいくつか書かなければならなかった後ではこの感情に共感できるが、
\newterm{型推論}と呼ばれる技術があって、コンパイラが使用されるコンテキストからほとん
どの型を推論できる。C++では、変数を\code{auto}と宣言してコンパイラにその型を決定させ
ることができる。

Haskellでは、まれな場合を除いて、型注釈は純粋にオプションだ。それでもプログラマーは
型注釈を使う傾向がある。なぜならコードの意味論について多くのことを伝えられるし、コン
パイルエラーを理解しやすくするからだ。Haskellではプロジェクトを型の設計から始めるのが
一般的な慣習だ。\sloppy{後に、型注釈が実装を導き、コンパイラによって強制されたコメントになる。}

強い静的型付けはコードのテストをしない言い訳としてよく使われる。Haskellプログラマーが
「コンパイルが通ったなら、正しいに違いない」と言うのを聞くことがあるかもしれない。も
ちろん、型が正しいプログラムが正しい出力を生成するという意味で正しいという保証はない。
この無頓着な態度の結果として、いくつかの研究ではHaskellがコード品質において期待される
ほど他よりも大きく前に出ていなかった。商業的な環境では、バグを修正するプレッシャーは
ある品質レベルまでしか適用されないようで、それはソフトウェア開発の経済学とエンドユー
ザーの許容度と大いに関係があり、プログラミング言語や方法論とはほとんど関係がない。より
良い基準は、スケジュールが遅れたり、大幅に縮小された機能で納品されたプロジェクトの数
を測ることだろう。

単体テストが強い型付けの代替になるという論拠については、強く型付けされた言語での一般
的なリファクタリング慣習を考えてみよう：特定の関数の引数の型を変更する。強く型付けされ
た言語では、その関数の宣言を変更してビルドエラーをすべて修正すれば十分だ。弱く型付けさ
れた言語では、関数が今や異なるデータを期待するという事実を呼び出し元に伝播できない。
単体テストは不一致のいくつかを捕まえるかもしれないが、テストはほぼ常に決定論的プロセス
ではなく確率的プロセスだ。テストは証明の貧弱な代替物だ。

\section{型とは何か？}

型の最もシンプルな直感は、それらが値の集合だということだ。型\code{Bool}（Haskellでは
具体的な型は大文字で始まることを覚えておこう）は\code{True}と\code{False}の2要素集合
だ。型\code{Char}は\code{a}や\code{ą}のようなすべてのUnicode文字の集合だ。

集合は有限でも無限でもよい。\code{Char}のリストの同義語である\code{String}の型は、
無限集合の例だ。

\code{x}を\code{Integer}として宣言するとき：

\src{snippet01}
それが整数の集合の要素であると言っている。Haskellの\code{Integer}は無限集合であり、
任意精度算術に使用できる。C++の\code{int}と同様に、機械型に対応する有限集合\code{Int}
もある。

型と集合のこの同一視をやっかいにする微妙な点がいくつかある。循環定義を含む多相関数に
関する問題や、すべての集合の集合を持てないという事実の問題がある。しかし約束したとおり、
数学の細部にこだわらないことにする。素晴らしいのは、$\Set$と呼ばれる集合の圏があり、
それを使えばいいということだ。$\Set$では、対象は集合で射（矢印）は関数だ。

$\Set$は非常に特別な圏だ。なぜなら実際にその対象の内部を覗くことができ、そうすること
で多くの直感を得られるからだ。例えば、空集合には要素がないことを知っている。特別な1要
素集合があることも知っている。関数はある集合の要素を別の集合の要素にマップすることを知
っている。2つの要素を1つにマップできるが、1つの要素を2つにはできない。恒等関数は集合の
各要素をそれ自身にマップすることなども知っている。計画はこれらすべての情報を徐々に忘れ、
代わりにすべてのそれらの概念を純粋に圏論的な用語、つまり対象と射の用語で表現することだ。

理想の世界では、Haskellの型は集合でHaskellの関数は集合間の数学的関数だと言えるだろう。
ただ一つの小さな問題がある：数学的関数はコードを実行しない――ただ答えを知っている。
Haskellの関数は答えを計算しなければならない。答えが有限のステップ数で得られるなら問題
ない――その数がいかに大きくても。しかし再帰を含む計算があり、それらは終了しないかもし
れない。終了する関数と終了しない関数を区別することは決定不可能だ――有名な停止問題がある
――からHaskellから終了しない関数を単純に禁止することはできない。そのためコンピュータ科
学者は、見方によっては天才的なアイデア、あるいは大きなハックを思いついた：\newterm{ボトム}
と呼ばれ\code{\_|\_}またはUnicode $\bot$で表記される特別な値で、すべての型を拡張する
というものだ。この「値」は終了しない計算に対応する。だから次のように宣言された関数は：

\src{snippet02}
\code{True}、\code{False}、または\code{\_|\_}を返す可能性があり、最後のものは決して
終了しないことを意味する。

興味深いことに、ボトムを型システムの一部として受け入れると、すべての実行時エラーをボト
ムとして扱い、さらに関数が明示的にボトムを返すことを許可するのが便利になる。後者は通常
\code{undefined}という式を使って行われる：

\src{snippet03}
この定義は型チェックを通過する。なぜなら\code{undefined}はボトムに評価され、\code{Bool}
を含む任意の型のメンバーだからだ。さらに次のように書くこともできる：

\src{snippet04}
（\code{x}なしで）ボトムは型\code{Bool -> Bool}のメンバーでもあるからだ。

ボトムを返す可能性のある関数は部分関数と呼ばれ、すべての可能な引数に対して有効な結果を
返す全関数と対比される。

ボトムのために、Haskellの型と関数の圏が$\Set$ではなく$\Hask$と呼ばれるのを見ることに
なるだろう。理論的な観点からは、これは尽きることのない複雑さの源なので、ここで私は肉切
り包丁を使いこの推論の流れを断ち切ることにする。実用的な観点からは、終了しない関数とボ
トムを無視し、$\Hask$を真正の$\Set$として扱っても問題ない。\footnote{Nils Anders
  Danielsson, John Hughes, Patrik Jansson, Jeremy Gibbons, \href{http://www.cs.ox.ac.uk/jeremy.gibbons/publications/fast+loose.pdf}{
    Fast and Loose Reasoning is Morally Correct}. この論文はほとんどの文脈でボトムを
  無視することの正当化を提供する。}

\section{なぜ数学的モデルが必要か？}

プログラマーとして、あなたはプログラミング言語の構文と文法によく精通している。言語のこ
れらの側面は通常、言語仕様の冒頭で形式的な記法を使って記述される。しかし言語の意味、
すなわちセマンティクスは記述がはるかに難しく、より多くのページを必要とし、形式的である
ことはまれで、ほぼ完全ではない。だからこそ、言語の専門家たちの終わりなき議論や、言語標
準の細部の解釈に特化した書籍の一大産業が生まれているのだ。

言語のセマンティクスを記述するための形式的なツールはあるが、その複雑さのために、現実の
プログラミングの大物ではなく、簡略化された学術的な言語と共にほとんど使用される。
\newterm{操作的意味論}と呼ばれるそのようなツールの一つは、プログラム実行のメカニクスを
記述する。それは形式化された理想化されたインタープリタを定義する。C++のような産業用言語
のセマンティクスは、通常「抽象機械」という観点から非形式的な操作的推論を使って記述される。

問題は、操作的意味論を使ってプログラムに関することを証明するのが非常に難しいということ
だ。プログラムの性質を示すには、基本的にそれを理想化されたインタープリタで「実行する」
必要がある。

プログラマーが正しさの形式的証明を決して行わないことは重要ではない。私たちはいつも正し
いプログラムを書いていると「考えている」。「ああ、とりあえずコードを数行投げてみて何が
起こるか見てみよう」などとキーボードの前で言う人はいない。私たちが書くコードが特定のア
クションを実行して望ましい結果をもたらすと思っている。それがそうでないとき、たいてい驚
く。それは、私たちが書くプログラムについて推論しており、通常は頭の中でインタープリタを
実行してそれを行っていることを意味する。ただすべての変数を追跡するのが本当に難しい。コン
ピュータはプログラムの実行が得意だが、人間はそうではない！もしそうなら、コンピュータは
必要ない。

しかし代替手段がある。\newterm{表示的意味論}と呼ばれ、数学に基づいている。表示的意味論
では、すべてのプログラムの構成物に数学的解釈が与えられる。それにより、プログラムの性質
を証明したければ、数学的定理を証明するだけでよい。定理の証明は難しいと思うかもしれない
が、実際のところ、人間は何千年も数学的手法を積み上げてきたので、活用できる蓄積された知
識が豊富にある。また、プロの数学者が証明する定理の種類と比べると、プログラミングで遭遇
する問題は通常かなり単純で、些末なことも多い。

表示的意味論にかなり適した言語であるHaskellでの階乗関数の定義を考えよう：

\src{snippet05}
式\code{{[}1..n{]}}は\code{1}から\code{n}までの整数のリストだ。関数\code{product}は
リストのすべての要素を掛け合わせる。それはまるで数学のテキストから取られた階乗の定義の
ようだ。これをCと比較しよう：

\begin{snip}{c}
int fact(int n) {
    int i;
    int result = 1;
    for (i = 2; i <= n; ++i)
        result *= i;
    return result;
}
\end{snip}
もっと言う必要があるだろうか？

わかった、これはちょっとずるい攻め方だったと最初に認めよう！階乗関数には明らかな数学
的表示がある。洞察力のある読者はこう問うかもしれない：キーボードから文字を読み込んだり、
ネットワーク越しにパケットを送信したりするための数学的モデルは何か？長い間、これは厄介
な質問で、かなり込み入った説明につながっていた。有用なプログラムを書くために不可欠な重
要なタスクのかなりの数について、表示的意味論は最良の適合ではないように思われ、それらは
操作的意味論で簡単に取り扱えた。突破口は圏論からやってきた。Eugenio Moggiが計算効果を
モナドにマップできることを発見した。これは重要な観察であり、表示的意味論に新たな命を吹
き込んで純粋関数型プログラムをより使いやすくしただけでなく、伝統的なプログラミングに新
たな光を当てた。より多くの圏論的ツールを開発したとき、後でモナドについて話す。

プログラミングの数学的モデルを持つことの重要な利点の一つは、ソフトウェアの正しさの形式
的証明を実行できることだ。コンシューマーソフトウェアを書いているとき、これはそれほど重
要に思えないかもしれないが、失敗のコストが法外だったり人命が危機にさらされているプログ
ラミングの分野もある。しかし医療システムのウェブアプリケーションを書くときでも、Haskell
標準ライブラリの関数とアルゴリズムが正しさの証明付きで提供されているという事実をありが
たく思うかもしれない。

\section{純粋な関数と汚い関数}

C++やその他の命令型言語で関数と呼ぶものは、数学者が関数と呼ぶものと同じではない。数学
的関数は単に値から値へのマッピングだ。

数学的関数はプログラミング言語で実装できる：そのような関数は入力値が与えられると出力値
を計算する。数の二乗を生成する関数はおそらく入力値をそれ自身に掛けるだろう。呼ばれるた
びにそれを行い、同じ入力で呼ばれるたびに同じ出力を生成することが保証されている。数の二
乗は月の満ち欠けで変わらない。

また、数の二乗を計算することが犬においしいおやつを与えるという副作用を持つべきではない。
そのような「関数」は数学的関数として簡単にモデル化できない。

プログラミング言語では、同じ入力に対して常に同じ結果を生成し副作用を持たない関数は
\newterm{純粋関数}と呼ばれる。Haskellのような純粋関数型言語では、すべての関数が純粋だ。
そのため、これらの言語に表示的意味論を与え、圏論を使ってモデル化するのが容易だ。他の言
語については、純粋なサブセットに制限したり、副作用について別々に推論することは常に可能
だ。後でモナドがどのように純粋な関数だけを使ってあらゆる種類の効果をモデル化できるかを
見ることになる。したがって、数学的関数に制限しても実際には何も失わない。

\section{型の例}

型が集合であることに気づくと、かなりエキゾチックな型について考えることができる。例えば、
空集合に対応する型は何か？いや、C++の\code{void}ではない。ただしこの型はHaskellでは
\code{Void}と呼ばれる。これはどんな値も住んでいない型だ。\code{Void}を取る関数を定義
できるが、決して呼び出せない。呼び出すには型\code{Void}の値を提供しなければならないが、
そんなものは存在しない。この関数が返せるものについては、何の制限もない。任意の型を返す
ことができる（呼び出せないので決して返さないが）。言い換えれば、戻り型において多相的な
関数だ。Haskellプログラマーはこれに名前を付けている：

\src{snippet06}
（覚えておこう、\code{a}は任意の型を表す型変数だ。）この名前は偶然ではない。カリー＝
ハワード同型と呼ばれる論理の観点から型と関数のより深い解釈がある。型\code{Void}は偽を
表し、関数\code{absurd}の型は偽からは何でも導かれるという命題に対応する。ラテン語の格言
「ex falso sequitur quodlibet（偽から何でも導かれる）」のように。

次は単集合に対応する型だ。可能な値が1つだけの型だ。この値はただ「ある」。すぐにはそう
認識しないかもしれないが、それがC++の\code{void}だ。この型から・への関数を考えよう。
\code{void}からの関数は常に呼び出せる。純粋関数であれば、常に同じ結果を返す。そのよう
な関数の例を示す：

\begin{snip}{c}
int f44() { return 44; }
\end{snip}
この関数は「何も」取らないと考えるかもしれないが、今見たように、「何も」取る関数は
「何も」を表す値が存在しないため決して呼び出せない。では、この関数は何を取るのか？概念
的には、インスタンスが1つしかないダミーの値を取るので、明示的に言及する必要はない。しか
しHaskellには、この値のシンボルがある：空の括弧のペア\code{()}だ。だから、面白い偶然
（本当に偶然だろうか？）で、voidの関数への呼び出しはC++とHaskellで同じに見える。また、
Haskellの簡潔さへの愛から、同じシンボル\code{()}が型、コンストラクタ、そして単集合に対
応する唯一の値に使われる。だからHaskellでのこの関数はこうだ：

\src{snippet07}
最初の行は\code{f44}が型\code{()}（「ユニット」と発音する）を型\code{Integer}にマップ
することを宣言している。2行目は\code{f44}をユニットの唯一のコンストラクタ\code{()}に
対するパターンマッチングで定義し、数44を生成する。この関数はユニット値\code{()}を提供
することで呼び出す：

\begin{snip}{c}
f44 ()
\end{snip}
ユニットのすべての関数は、目標の型から単一の要素を選ぶことと同等であることに注意しよう
（ここでは整数44を選ぶ）。実際、\code{f44}を数44の別の表現として考えることができる。
これは集合の要素の明示的な言及を、代わりに関数（射）について話すことで置き換えられる方
法の例だ。ユニットから任意の型$A$への関数は、集合$A$の要素と一対一対応する。

戻り型がvoidの関数、あるいはHaskellでユニット戻り型の関数はどうだろう？C++ではそのよう
な関数は副作用のために使われるが、これらは数学的な意味での本当の関数ではないことを知って
いる。ユニットを返す純粋関数は何もしない：引数を捨てる。

数学的には、集合$A$から単集合への関数は$A$のすべての要素をその単集合の単一要素にマップ
する。すべての$A$に対して、そのような関数がちょうど1つ存在する。\code{Integer}のための
この関数はこうだ：

\src{snippet08}
任意の整数を与えると、ユニットを返す。簡潔さの精神から、Haskellでは破棄される引数にワイ
ルドカードパターン、つまりアンダースコアを使用できる。こうすることで名前を考える必要がな
い。したがって、上記は次のように書き直せる：

\src{snippet09}
この関数の実装は渡された値に依存しないだけでなく、引数の型にさえ依存しないことに注意し
よう。

任意の型に対して同じ公式で実装できる関数は、パラメトリック多相的と呼ばれる。具体的な型
の代わりに型パラメータを使って1つの方程式でそのような関数のファミリー全体を実装できる。
任意の型からユニット型へのこの多相関数を何と呼ぶべきか？もちろん\code{unit}と呼ぼう：

\src{snippet10}
C++ではこの関数をこう書くだろう：

\begin{snip}{cpp}
template<typename T>
void unit(T) {}
\end{snip}
型の類型において次は2要素集合だ。C++では\code{bool}と呼ばれ、Haskellでは予想通り
\code{Bool}と呼ばれる。違いはC++の\code{bool}が組み込み型であるのに対し、Haskellでは
次のように定義できる：

\src{snippet11}
（この定義の読み方は、\code{Bool}は\code{True}か\code{False}のどちらかだということだ。）
原則として、C++でも列挙型としてブール型を定義できるはずだ：

\begin{snip}{cpp}
enum bool {
    true,
    false
};
\end{snip}
しかしC++の\code{enum}は内部的に整数だ。代わりにC++11の「\code{enum class}」を使うこと
もできるが、その場合\code{bool::true}や\code{bool::false}のようにクラス名で値を修飾し
なければならないし、それを使用するすべてのファイルで適切なヘッダをインクルードしなければ
ならない。

\code{Bool}からの純粋関数は、目標の型から2つの値を選ぶだけだ。1つは\code{True}に対応し、
もう1つは\code{False}に対応する。

\code{Bool}への関数は\newterm{述語}と呼ばれる。例えば、HaskellライブラリのData.Charは
\code{isAlpha}や\code{isDigit}のような述語で満ちている。C++には同様のライブラリ
\code{<cctype>}があり、\code{isalpha}や\code{isdigit}などを定義しているが、これらはブール
値ではなく\code{int}を返す。実際の述語は\code{std::ctype}に定義されており、
\code{ctype::is(alpha, c)}、\code{ctype::is(digit, c)}などの形式を持つ。

\section{課題}

\begin{enumerate}
  \tightlist
  \item
        好きな言語で高階関数（または関数オブジェクト）\code{memoize}を定義せよ。この
        関数は純粋関数\code{f}を引数に取り、\code{f}とほぼ同じように振る舞うが、すべての
        引数に対して元の関数を一度だけ呼び出し、結果を内部に格納し、同じ引数で呼ばれる
        たびにこの格納された結果を返す関数を返す。メモ化された関数と元の関数の違いはパ
        フォーマンスで確認できる。例えば、評価に時間がかかる関数をメモ化してみよう。最初
        に呼び出すときは結果を待たなければならないが、同じ引数での2回目以降の呼び出しでは、
        すぐに結果が得られるはずだ。
  \item
        標準ライブラリから乱数を生成するために通常使う関数をメモ化してみよう。うまくいくか？
  \item
        ほとんどの乱数生成器はシードで初期化できる。シードを取り、そのシードで乱数生成器
        を呼び出し、結果を返す関数を実装せよ。その関数をメモ化せよ。うまくいくか？
  \item
        これらのC++関数のうち、どれが純粋か？それらをメモ化し、メモ化した場合としない場合
        で複数回呼び出すと何が起こるかを観察せよ。

        \begin{enumerate}
          \tightlist
          \item
                本文の例にある階乗関数。
          \item
                \begin{minted}{cpp}
std::getchar()
\end{minted}
          \item
                \begin{minted}{cpp}
bool f() {
    std::cout << "Hello!" << std::endl;
    return true;
}
\end{minted}
          \item
                \begin{minted}{cpp}
int f(int x) {
    static int y = 0;
    y += x;
    return y;
}
\end{minted}
        \end{enumerate}
  \item
        \code{Bool}から\code{Bool}への異なる関数はいくつあるか？すべて実装できるか？
  \item
        対象が型\code{Void}、\code{()}（ユニット）、\code{Bool}だけからなる圏の絵を描け。
        矢印はこれらの型間の可能なすべての関数に対応するものとし、矢印に関数の名前の
        ラベルを付けよ。
\end{enumerate}
